\documentclass[11pt]{article}
\usepackage[margin=2.1cm]{geometry}
\usepackage{amsmath,amssymb}
\usepackage{enumitem}
\usepackage{parskip}
\usepackage{booktabs}
\usepackage{xcolor}
\usepackage[hidelinks]{hyperref}

\setlist{nosep, leftmargin=1.4em}
\newcommand{\solid}{\textcolor{teal}{\textbf{[solid]}}}
\newcommand{\guess}{\textcolor{orange!80!black}{\textbf{[conjecture]}}}

\begin{document}

{\Large\textbf{Meeting prep: expiry-day microstructure and the settlement VWAP}}\\
For my own reading before meeting Prof.\ Banerjee. What we did, what we found, what we only
think, and what he may ask.

\hrule
\vspace{0.8em}

\section*{0. The story in five lines}
\begin{enumerate}
  \item On expiry days more trading is drawn into the 15:00 to 15:30 settlement window of F\&O
  stocks, and that window becomes more volatile and more expensive to trade in.
  \item Nobody pushes \emph{stock} settlement prices: stock F\&O settle by delivery, so there is
  nothing to gain. Order flow in the window stays two-sided.
  \item The \emph{index} (cash-settled) is the exception: Nifty 50 falls 6.2 bps against Nifty
  Next 50 over the window on index expiries.
  \item The next-morning fall in Figure 2 is mostly information arriving overnight, not the
  settlement price being released.
  \item In the expiring stock future the spread widens three times as much as in cash in the
  last half hour, because market makers leave a contract about to turn into delivery.
\end{enumerate}

\solid{} = backed by a test. \guess{} = our explanation, not yet tested directly.

\section*{1. What we did, step by step}
\begin{enumerate}
  \item \textbf{Data.} NSE order and trade files for 2022. Cash market: trades for 242 sessions
  (one corrupt file, 2 Aug); orders for 24 sessions (12 monthly expiries + the day before each).
  Index file: Nifty 50 and Nifty Next 50 every second. F\&O trades and orders: all of 2022, from
  the library's shares (we first had only Mar to Jun).
  \item \textbf{Parsing.} Our own Rust parser (nsetick) turns the fixed-width NSE files into
  parquet, split by date and stock.
  \item \textbf{Universe.} 150 stocks: 50 liquid F\&O, 50 illiquid F\&O, 50 \emph{placebo} stocks
  with no derivatives, matched on turnover and price. Placebo stocks trade on the same days but
  have nothing settling, so anything also seen there is about the day, not the settlement.
  \item \textbf{Trade direction.} Every trade has the buy and sell order numbers; order numbers
  are given in order of arrival, so the later order is the one that crossed the spread (the
  aggressor). So we know who initiated every trade exactly.
  \item \textbf{Order book.} For the 24 sessions we replay every order entry, change and
  cancellation to rebuild the full limit order book, and take a snapshot every 5 seconds.
  Gives quoted spread, depth, cancellations, hidden size, etc.
  \item \textbf{Daily tape.} For all 242 sessions, one row per stock per day: settlement VWAP,
  price before the window, VWAP of each 5-minute block, buyer- and seller-initiated volume,
  participant categories, next morning's prices. 36,300 stock-days.
  \item \textbf{Tests.} Paired tests expiry vs.\ day before (order book sample); regressions on
  the full year with stock and session fixed effects; placebo comparisons; randomization
  inference.
  \item \textbf{After the 21 Oct meeting.} News tests for Figure 2; spread and depth through
  the whole day; volatility regressions; futures trades and futures order books for all 12
  expiries.
\end{enumerate}

\section*{2. Methods in plain words (in case he asks)}
\begin{itemize}
  \item \textbf{Expiry vs.\ day before (paired).} Same stock, same month, one day apart. Removes
  anything specific to the stock or the month.
  \item \textbf{Two-way fixed effects.} Regression with a dummy for every stock and every
  session. The session dummy removes anything that hit the whole market that day (e.g.\ a crash),
  so a market-wide move cannot look like an expiry effect. Standard errors clustered by session
  (stocks on the same day are not independent).
  \item \textbf{Placebo / difference-in-differences.} Compare the change in F\&O stocks with the
  change in placebo stocks on the same days. If both change equally, it is the day, not the
  settlement.
  \item \textbf{Randomization inference.} Pretend a random ordinary session in each month is the
  ``expiry'', re-run the regression 1,000 times; the $p$-value is how often a fake expiry gives
  an effect as big as the real one. Does not rely on normal-distribution assumptions.
  \item \textbf{Four broad hypotheses, one primary measure each} (as he asked), with false
  discovery rate control across the four primaries; everything else is supporting.
  \item \textbf{OIB} $= (BI - SI)/(BI + SI)$, buyer- minus seller-initiated volume over total.
  $|$OIB$|$ near 0 means two-sided, near 1 means one-sided.
  \item \textbf{Realised variance (volatility)}: sum of squared 5-second changes of the mid
  price. Our measure of how much prices move.
  \item \textbf{Log-ratio \%}: $100 \times \log(\text{expiry}/\text{day before})$, roughly the
  percentage change.
\end{itemize}

\section*{3. Results by hypothesis}

\subsection*{BH1: Is the settlement price pushed?}
\textbf{Primary:} distance of the settlement VWAP from the mid at 15:00. 36.0 bps on expiry vs.\
34.5 on the day before, $p=0.21$ (1,792 pairs). \textbf{Not rejected.} \solid

Supporting:
\begin{itemize}
  \item Window share of daily volume rises on monthly expiries, 2.4 pp more in F\&O stocks than
  placebo ($p<0.001$). More trading, but the price does not move further. \solid
  \item Futures pricing: in the last minutes the expiring future moves with the live cash price,
  slope 0.85 (se 0.04), not with the VWAP (which by 15:28 is 98\% fixed). That is what physical
  delivery implies. \solid
  \item No pull towards option strikes: distance to nearest strike unchanged. \solid
  \item \textbf{Index:} Nifty 50 minus Nifty Next 50 falls 6.18 bps more over the window on
  index expiries (about 11 index points), randomization $p=0.008$; negative on 63\% of index
  expiries vs.\ 52\% of other days; same on weekly and monthly. \solid
\end{itemize}
\textbf{Why:} stock F\&O are physically settled, so moving the VWAP gains nothing (a long future
pays the futures price for the shares whatever the settlement). The index is cash-settled, so
there the incentive is real. \solid{} for the stock part (futures slope).

\textbf{Why does the index fall?} \guess
\begin{enumerate}
  \item Hedge unwinding: arbitrageurs hold the stocks and are short futures/options; they exit
  by selling the stocks in the window at the settlement average. Selling pressure in Nifty 50
  stocks only.
  \item Deliberate pressure by large index-option holders.
\end{enumerate}
We cannot tell them apart without positions. The index recovers about the same amount next
morning (6.8 bps) but imprecisely ($p=0.28$), which fits temporary selling pressure. Next test:
aggressor flow by participant type in Nifty 50 stocks on index expiries.

\subsection*{BH2: Is the window costlier to trade in?}
\textbf{Primary:} quoted spread. 5.31 $\to$ 5.60 bps, $p<0.001$; specific to F\&O stocks (vs.\
placebo $+0.28$ bps, $p=0.016$). \textbf{Rejected.} \solid

New plots (full day and minute by minute):
\begin{center}
\small
\begin{tabular}{lrrr}
\toprule
Window, expiry vs.\ day before & Spread & Depth & Volatility \\
\midrule
Liquid F\&O & $+9.8$\% & $-7.2$\% & $+45$\% \\
Illiquid F\&O & $+10.3$\% & $-4.8$\% & $+46$\% \\
Placebo & $+1.2$\% n.s. & $-0.5$\% n.s. & $+7$\% n.s. \\
\bottomrule
\end{tabular}
\end{center}
Before the window the spread is only slightly wider (liquid $+4.5$\%), depth unchanged.

\textbf{Why: volatility.} \solid{} Regress each stock-month's change in spread on its change in
volatility. Volatility explains about 70\% of the widening and all of the fall in depth.
Market makers quote wider and smaller when prices move more (inventory risk, risk of trading
with someone better informed).

\textbf{Why is the window more volatile on expiry?} \guess
\begin{itemize}
  \item More trading pulled into the window (closing, rolling, benchmarking to the close).
  \item Option hedging: on expiry day options near the strike switch quickly between worthless
  and in the money (very high gamma), so dealers hedging in cash must trade a lot in the last
  half hour.
  \item The residual 3\% widening: market makers also active in the expiring options may be busy
  there.
\end{itemize}

\subsection*{BH3: Is there one-sided pressure?}
\textbf{Primary:} $|$OIB$|$ over the window. 0.12 on ordinary days (liquid), $+0.02$ on expiries
(n.s.); paired $p=0.54$. \textbf{Not rejected.} \solid

Supporting:
\begin{itemize}
  \item OIB is more \emph{persistent} block to block on expiries (0.34 $+$ 0.15): large orders
  worked through the window. \solid
  \item Variance ratio and permanent share of price impact unchanged (0.75 vs.\ 0.76; 1.19 vs.\
  1.22): no sign of sustained pressure. \solid
  \item Proprietary traders trade \emph{against} the window move and earn it back overnight
  (liquidity provision, not pushing). Not different on expiries. \solid
  \item Custodian (institutional) share of window volume 59\% $\to$ 65\%. \solid
  \item Futures: the expiring contract's $|$OIB$|$ in the window unchanged (0.18); roll mostly
  done before (next month has 52\% of futures volume the day before, 59\% on expiry). \solid
\end{itemize}
\textbf{Why:} no gain from pushing; closing flow comes from both sides; positions already
rolled. Persistence = institutions slicing orders evenly to get the closing price; across stocks
and institutions buys and sells offset. \guess{} for the slicing interpretation.

\subsection*{BH4: Does order behaviour change?}
\textbf{Primary:} orders cancelled within 1 second, 28.6\% $\to$ 30.1\%, $p<0.001$; 2.4 pp more
than placebo. \textbf{Rejected.} \solid

Supporting: entries rise as fast as cancellations (more quoting overall, not place-and-pull);
hidden share of the book high (82\%) and unchanged; IOC use unchanged; proprietary volume up.

\textbf{Why: volatility again.} \solid{} Controlling for each stock's rise in volatility, the
increase is small and not significant. Algorithms re-price every time the price moves.

\section*{4. Figure 2: is the next-morning fall news?}
Figure 2 = average path through the window and to the next morning, oriented so the window
move is positive. On monthly expiries part of the move is given back next morning.
Regression: extra reversal $-0.26$ ($p=0.04$) in F\&O stocks, $-0.23$ in placebo.

What we did:
\begin{itemize}
  \item Dropped stock-days with signs of news (next-day volume $> 2\times$ median, or opening gap
  in the stock's top 5\%).
  \item Adjusted for market news: stock beta $\times$ market's overnight move.
  \item Dropped 24 Feb 2022 (Ukraine invasion: crash on the day, rebound next morning).
  \item Estimated expiry by expiry.
\end{itemize}
\begin{center}
\small
\begin{tabular}{lcc}
\toprule
Extra reversal on expiries & F\&O & Placebo \\
\midrule
Baseline & $-0.26$ ($p=0.04$) & $-0.23$ ($p=0.007$) \\
Beta-adjusted, no news days & $-0.13$ ($p=0.34$) & $-0.06$ ($p=0.17$) \\
\bottomrule
\end{tabular}
\end{center}
\textbf{Conclusion:} mostly news. \solid{} Comes from a few expiries (24 Feb, 31 Mar, 25 Aug,
29 Sep); 10 of 12 inside the normal range; as strong in placebo; insignificant after
adjustments.

\textbf{How news makes a ``reversal''} (if he asks why news should matter at all): unexpected
news is unrelated to the window's direction, so on average it cannot undo it. But with 12
expiries a few big days dominate. On 24 Feb high-beta stocks fell most in the window and bounced
most next morning: looks like a reversal stock by stock. Beta adjustment removes half.

\textbf{Caveat / residual:} the morning-to-15:00 move is reversed much less ($-0.05$) than the
window move, so maybe not all news. \guess: closing orders from funds near month end absorbed by
market makers, who earn a small reversal next morning (would hit placebo too). Test: company
announcement data.

\section*{5. Futures (all 12 expiries)}
Rebuilt the order book of each expiring stock future (84 to 100 stocks per day).
\begin{center}
\small
\begin{tabular}{lrrr}
\toprule
Window, expiry vs.\ day before & Spread & Depth & Volatility \\
\midrule
Cash (same stocks) & $+10.0$\% & $-6.0$\% & $+45.8$\% \\
Expiring future & $+33.4$\% & $-6.4$\% & $+58.7$\% \\
\bottomrule
\end{tabular}
\end{center}
\begin{itemize}
  \item Before the window the future looks normal (spread not wider). The jump is only in the
  last half hour. \solid
  \item Volatility explains only a third of it (vs.\ 70\% in cash); depth fall significant even
  after controlling for volatility. \solid
  \item \guess: market makers normally carry a small position overnight and offset it next
  morning. On expiry there is no next morning: anything open at 15:30 becomes a delivery of
  shares (needs shares or cash, extra margin). So they pull out of the expiring contract;
  liquidity has moved to next month anyway.
  \item Still no one-sided pressure in futures. \solid
\end{itemize}

\section*{6. Solid vs.\ conjecture at a glance}
\begin{center}
\small
\begin{tabular}{p{7.2cm}p{8.4cm}}
\toprule
Solid (tested) & Conjecture (our explanation) \\
\midrule
No push in stock settlement prices & Index fall = hedge unwinding vs.\ deliberate push \\
Futures priced on cash price (delivery) & Options too small per stock to be worth pushing \\
Index falls 6.2 bps on index expiries & Window volatility from option (gamma) hedging \\
Window costlier, only F\&O stocks & Residual 3\% from market makers busy in options \\
Volatility explains most of spread/depth/fleeting & OIB persistence = institutions slicing orders \\
No one-sided OIB, in cash or futures & Residual overnight reversal = closing-order compensation \\
Next-morning fall mostly news & Futures widening = delivery risk for market makers \\
Futures window spread $+33$\% \\
\bottomrule
\end{tabular}
\end{center}

\section*{7. Questions he may ask, and answers}
\begin{itemize}
  \item \textbf{``So is there manipulation?''} Not in stock settlement prices; the incentive
  is not there (delivery). The index settlement is systematically lower, which is where the
  incentive is, but without positions or trader IDs we cannot say who or why.
  \item \textbf{``How do you know the aggressor?''} Order numbers are sequential by arrival; the
  later order crossed the spread. Exact, no Lee--Ready style guessing.
  \item \textbf{``Why the day before as control?''} Same stock, same month, same contract one
  day from expiry. We also use all other sessions in the full-year regressions.
  \item \textbf{``Why should news cause a reversal?''} See section 4: it should not on average;
  a few big news days plus beta differences created it.
  \item \textbf{``How do you know it's volatility and not something else?''} Cross-stock: stocks
  whose volatility rose more widened more; at zero volatility change the widening is only 3\%.
  Also placebo stocks: no extra volatility, no widening.
  \item \textbf{``But why is volatility higher?''} Honest answer: more trading in the window,
  and probably option hedging near expiry. That part is a hypothesis.
  \item \textbf{``Did you check futures?''} Yes, full order books for all 12 expiries. Huge
  widening in the last half hour of the expiring contract; no one-sided pressure.
  \item \textbf{``Orders match his files?''} We left that check for later (our orders and trades
  reconcile with each other).
  \item \textbf{``What next?''} Company announcements (news test); participant flow in Nifty 50
  stocks on index expiries (index hypothesis); open interest if available.
\end{itemize}

\section*{8. Limitations to mention if asked}
\begin{itemize}
  \item 12 expiries per year: power is limited, a few days can matter.
  \item No positions/open interest and no trader IDs, only participant categories.
  \item Index settlement value approximated by its window average (checked within a few points).
  \item News measured through its footprint (volume, gaps), not actual announcements.
\end{itemize}

\end{document}
